% Part: first-order-logic
% Chapter: sequent-calculus
% Section: proving-things

\documentclass[../../../include/open-logic-section]{subfiles}

\begin{document}

\iftag{FOL}
      {\olfileid{fol}{seq}{pro}}
      {\olfileid{pl}{seq}{pro}}

\olsection{Voorbeelden: zo bouw je \usetoken{P}{derivation}}

\begin{ex}
Maak een $\Log{LK}$-!!{derivation} van de sequent $!A \land !B \Sequent !A$.

We zetten eerst de gewenste eindsequent onderaan de !!{derivation}.
\begin{prooftree}
\AxiomC{}
\UnaryInf$!A\land !B \fCenter !A$
\end{prooftree}
Nu moeten we uitzoeken welke afleidingsstap deze onderste sequent kan
hebben opgeleverd. Dat kan een structurele regel zijn. Toch is het
handig om eerst naar een logische regel te kijken. In de onderste
sequent komt maar één logisch connectief voor: $\land$. We hebben dus
een $\land$-regel nodig. Het $\land$-teken staat in het antecedent,
oftewel aan de linkerkant van de sequent. Daarom nemen we de regel
\LeftR{\land}.
\begin{prooftree}
\AxiomC{}
\RightLabel{\LeftR{\land}}
\UnaryInf$!A\land !B \fCenter !A$
\end{prooftree}
Boven deze toepassing van \LeftR{\land} kan een van twee sequenten
staan: $!A \Sequent !A$ of $!B \Sequent !A$. De eerste sequent, $!A
\Sequent !A$, is een beginsequent; dat is precies wat we zoeken. De
tweede, $!B \Sequent !A$, is in het algemeen niet afleidbaar. We vullen
daarom de eerste sequent in:
\begin{prooftree}
\Axiom$!A \fCenter !A$
\RightLabel{\LeftR{\land}}
\UnaryInf$!A\land !B \fCenter !A$
\end{prooftree}
Daarmee hebben we een correcte $\Log{LK}$-!!{derivation} van de
sequent $!A \land !B \Sequent !A$.
\end{ex}

\begin{ex}
Maak een $\Log{LK}$-!!{derivation} van de sequent $\lnot !A \lor !B
\Sequent !A \lif !B$.

Begin weer onderaan de !!{derivation} en schrijf daar de gewenste
eindsequent.
\begin{prooftree}
\AxiomC{}
\UnaryInf$\lnot !A \lor !B \fCenter !A \lif !B$
\end{prooftree}
Om te zien welke logische regel als laatste kan zijn gebruikt,
bekijken we de logische tekens in de eindsequent: $\lnot$, $\lor$ en
$\lif$. Op dit moment tellen alleen $\lor$ en $\lif$. Dat zijn de
!!{main operator}s van !!{sentence}s in de eindsequent. Het
$\lnot$-teken staat binnen het bereik van een ander connectief; dat
handelen we later af. Voor de laatste stap kunnen we dus \LeftR{\lor}
of \RightR{\lif} kiezen. Allebei kan. We nemen \RightR{\lif}, want dan
hoeven we de afleiding nog niet meteen in twee takken te splitsen. Om
de gegeven onderste sequent met \RightR{\lif} te krijgen, moet de
afleiding er zo uitzien:
\begin{prooftree}
\AxiomC{}
\UnaryInf$ !A, \lnot !A \lor !B \fCenter !B $
\RightLabel{\RightR{\lif}} \UnaryInf$ \lnot !A \lor !B \fCenter !A \lif !B $
\end{prooftree}

We kunnen de regel \LeftR{\lor} toepassen zodra $\lnot !A \lor !B$
helemaal links in het antecedent staat. Volgens het regelschema splitst
de afleiding dan in de volgende twee bovenste sequenten:
\begin{prooftree}
\AxiomC{}
\UnaryInf$\lnot !A, !A \fCenter !B$
\AxiomC{}
\UnaryInf$!B, !A \fCenter !B$
\RightLabel{\LeftR{\lor}}
\BinaryInf$ \lnot !A \lor !B, !A \fCenter !B $
\RightLabel{\RightR{\Exchange}}
\UnaryInf$ !A, \lnot !A \lor !B \fCenter !B $
\RightLabel{\RightR{\lif}}
\UnaryInf$ \lnot !A \lor !B \fCenter !A \lif !B $
\end{prooftree}
We werken van onder naar boven en willen bij beginsequenten uitkomen.
De rechtertak is al bijna klaar: na één verzwakking en één verwisseling
staat er bovenaan een beginsequent.
\begin{prooftree}
\AxiomC{}
\UnaryInf$\lnot !A, !A \fCenter !B$
\Axiom$!B \fCenter !B$
\RightLabel{\LeftR{\Weakening}}
\UnaryInf$!A, !B \fCenter !B$
\RightLabel{\LeftR{\Exchange}}
\UnaryInf$!B, !A \fCenter !B$
\RightLabel{\LeftR{\lor}}
\BinaryInf$\lnot !A \lor !B, !A \fCenter !B $
\RightLabel{\RightR{\Exchange}}
\UnaryInf$ !A, \lnot !A \lor !B \fCenter !B $
\RightLabel{\RightR{\lif}}
\UnaryInf$ \lnot !A \lor !B \fCenter !A \lif !B $
\end{prooftree}

Kijk nu naar de linkertak. Het enige logische connectief in
!!{sentence} is daar $\lnot$. De !!{sentence}s met dat teken staan in
het antecedent. Daarom gebruiken we de regel \LeftR{\lnot}.
\begin{prooftree}
\AxiomC{}
\UnaryInf$ !A \fCenter !B, !A$
\RightLabel{\LeftR{\lnot}}
\UnaryInf$\lnot !A, !A \fCenter !B$
\Axiom$!B \fCenter !B$
\RightLabel{\LeftR{\Weakening}}
\UnaryInf$!A, !B \fCenter !B$
\RightLabel{\LeftR{\Exchange}}
\UnaryInf$!B, !A \fCenter !B$
\RightLabel{\LeftR{\lor}}
\BinaryInf$\lnot !A \lor !B, !A \fCenter !B $
\RightLabel{\RightR{\Exchange}}
\UnaryInf$ !A, \lnot !A \lor !B \fCenter !B $
\RightLabel{\RightR{\lif}}
\UnaryInf$ \lnot !A \lor !B \fCenter !A \lif !B $
\end{prooftree}
Ook de linkertak is dan bijna klaar. Net als rechts zijn nog één
verzwakking en één verwisseling nodig om op een beginsequent uit te
komen.
\begin{prooftree}
\Axiom$!A \fCenter !A$
\RightLabel{\RightR{\Weakening}}
\UnaryInf$ !A \fCenter !A, !B$
\RightLabel{\RightR{\Exchange}}
\UnaryInf$ !A \fCenter !B, !A$
\RightLabel{\LeftR{\lnot}}
\UnaryInf$\lnot !A, !A \fCenter !B$
\Axiom$!B \fCenter !B$
\RightLabel{\LeftR{\Weakening}}
\UnaryInf$!A, !B \fCenter !B$
\RightLabel{\LeftR{\Exchange}}
\UnaryInf$!B, !A \fCenter !B$
\RightLabel{\LeftR{\lor}}
\BinaryInf$\lnot !A \lor !B, !A \fCenter !B $
\RightLabel{\RightR{\Exchange}}
\UnaryInf$ !A, \lnot !A \lor !B \fCenter !B $
\RightLabel{\RightR{\lif}}
\UnaryInf$ \lnot !A \lor !B \fCenter !A \lif !B $
\end{prooftree}
\end{ex}

\begin{ex}
Maak een $\Log{LK}$-!!{derivation} van de sequent $\lnot !A \lor \lnot !B
\Sequent \lnot (!A \land !B)$.

We gebruiken dezelfde aanpak als hierboven en schrijven eerst de
gewenste eindsequent onderaan.
\begin{prooftree}
\AxiomC{}
\UnaryInf$ \lnot !A \lor \lnot !B \fCenter \lnot (!A \land !B) $
\end{prooftree}
De hoofdconnectieven van de !!{sentence}s in de eindsequent zijn
$\lor$ en $\lnot$. We kunnen hier met \LeftR{\lor} of met
\RightR{\lnot} beginnen. We nemen \RightR{\lnot}; dan hoeven we de
afleiding nog even niet in twee takken te splitsen.
\begin{prooftree}
\AxiomC{}
\UnaryInf$!A \land !B, \lnot !A \lor \lnot !B \fCenter $
\RightLabel{\RightR{\lnot}}
\UnaryInf$\lnot !A \lor \lnot !B \fCenter \lnot (!A \land !B)$
\end{prooftree}
Nu kunnen we verder met \LeftR{\land} of met \LeftR{\lor}. We proberen
eerst \LeftR{\land}. Daarboven kan dan $!A, \lnot !A \lor !B \Sequent
\quad$ of $!B, \lnot !A \lor !B \Sequent \quad$ staan. De
!!{derivation} werkt voor $!A$ en $!B$ op dezelfde manier, dus kiezen
we de eerste sequent:
\begin{prooftree}
\AxiomC{}
\UnaryInf$!A, \lnot !A \lor \lnot !B \fCenter $
\RightLabel{\LeftR{\land}}
\UnaryInf$!A \land !B, \lnot !A \lor \lnot !B \fCenter $
\RightLabel{\RightR{\lnot}}
\UnaryInf$\lnot !A \lor \lnot !B \fCenter \lnot (!A \land !B)$
\end{prooftree}
Als we de !!{derivation} verder invullen, zien we het probleem:
\begin{prooftree}
\Axiom$!A \fCenter !A$
\RightLabel{\LeftR{\lnot}}
\UnaryInf$ \lnot !A, !A \fCenter$
\AxiomC{}
\RightLabel{?}
\UnaryInf$!A \fCenter !B$
\RightLabel{\LeftR{\lnot}}
\UnaryInf$ \lnot !B, !A \fCenter$
\RightLabel{\LeftR{\lor}}
\BinaryInf$\lnot !A \lor \lnot !B, !A \fCenter $
\RightLabel{\LeftR{\Exchange}}
\UnaryInf$!A, \lnot !A \lor \lnot !B \fCenter $
\RightLabel{\LeftR{\land}}
\UnaryInf$!A \land !B, \lnot !A \lor \lnot !B \fCenter $
\RightLabel{\RightR{\lnot}}
\UnaryInf$\lnot !A \lor \lnot !B \fCenter \lnot (!A \land !B)$
\end{prooftree}
De sequent bovenaan de rechtertak kan niet verder worden aangepakt.
Ook met structurele regels kunnen we er geen beginsequent van maken.
Dit pad loopt dus dood. De keuze voor \LeftR{\land} was verkeerd. We
gaan terug naar de vorige tussenstand en gebruiken daar \LeftR{\lor}
in plaats van die regel. Dan krijgen we
\begin{prooftree}
\AxiomC{}
\UnaryInf$\lnot !A, !A \land !B \fCenter $

\AxiomC{}
\UnaryInf$\lnot !B, !A \land !B \fCenter $

\RightLabel{\LeftR{\lor}}
\BinaryInf$\lnot !A \lor \lnot !B, !A \land !B \fCenter $
\RightLabel{\LeftR{\Exchange}}
\UnaryInf$!A \land !B, \lnot !A \lor \lnot !B \fCenter $
\RightLabel{\RightR{\lnot}}
\UnaryInf$\lnot !A \lor \lnot !B \fCenter \lnot (!A \land !B)$
\end{prooftree}
We maken beide takken af zoals we eerder hebben gedaan. Het resultaat
is
\begin{prooftree}
\Axiom$ !A \fCenter!A$
\RightLabel{\LeftR{\land}}
\UnaryInf$!A \land !B \fCenter !A$
\RightLabel{\LeftR{\lnot}}
\UnaryInf$\lnot !A, !A \land !B \fCenter $

\Axiom$ !B \fCenter !B$
\RightLabel{\LeftR{\land}}
\UnaryInf$!A \land !B \fCenter !B$
\RightLabel{\LeftR{\lnot}}
\UnaryInf$\lnot !B, !A \land !B \fCenter $

\RightLabel{\LeftR{\lor}}
\BinaryInf$\lnot !A \lor \lnot !B, !A \land !B \fCenter $
\RightLabel{\LeftR{\Exchange}}
\UnaryInf$!A \land !B, \lnot !A \lor \lnot !B \fCenter $
\RightLabel{\RightR{\lnot}}
\UnaryInf$\lnot !A \lor \lnot !B \fCenter \lnot (!A \land !B)$
\end{prooftree}
(We hadden bij deze stappen de $\land$-regels ook lager dan de
$\lnot$-regels kunnen zetten. Ook dan hadden we een correcte
!!{derivation} gekregen.)
\end{ex}

\begin{ex}
Tot nu toe hebben we de contractieregel niet gebruikt, maar soms
hebben we hem wel nodig. Stel dat we $\quad \Sequent !A \lor \lnot !A$
willen bewijzen. Als we $\RightR{\lor}$ van onder naar boven toepassen,
krijgen we een van deze twee !!{derivation}s:
\begin{prooftree}
\AxiomC{}
\UnaryInf$ \fCenter !A$
\RightLabel{\RightR{\lor}}
\UnaryInf$ \fCenter !A \lor \lnot !A$
\DisplayProof\qquad\bottomAlignProof
\AxiomC{}
\UnaryInf$!A \fCenter $
\RightLabel{\RightR{\lnot}}
\UnaryInf$ \fCenter \lnot !A$
\RightLabel{\RightR{\lor}}
\UnaryInf$ \fCenter !A \lor \lnot !A$
\end{prooftree}
Geen van beide komt bovenaan uit op een beginsequent. De oplossing is
om de contractieregel van onder naar boven te lezen. Van boven naar
beneden voegt die regel twee voorkomens van !!a{sentence} samen tot één
voorkomen. Als we van onder naar boven een bewijs zoeken, mogen we dus
in de premisse een extra voorkomen van $!A \lor \lnot !A$ laten staan.
Dat kan bijvoorbeeld zo:
\begin{prooftree}
\AxiomC{}
\UnaryInf$ \fCenter !A \lor \lnot !A, !A$
\RightLabel{\RightR{\lor}}
\UnaryInf$ \fCenter !A \lor \lnot !A, !A \lor \lnot !A$
\RightLabel{\RightR{\Contraction}}
\UnaryInf$ \fCenter !A \lor \lnot !A$
\end{prooftree}
Nu kunnen we $\RightR{\lor}$ nog een keer toepassen en ook $\lnot !A$
krijgen. Daarmee wordt de !!{derivation} volledig.
\begin{prooftree}
\Axiom$!A \fCenter !A$
\RightLabel{\RightR{\lnot}}
\UnaryInf$\fCenter !A, \lnot !A$
\RightLabel{\RightR{\lor}}
\UnaryInf$\fCenter !A, !A \lor \lnot !A$
\RightLabel{\RightR{\Exchange}}
\UnaryInf$ \fCenter !A \lor \lnot !A, !A$
\RightLabel{\RightR{\lor}}
\UnaryInf$ \fCenter !A \lor \lnot !A, !A \lor \lnot !A$
\RightLabel{\RightR{\Contraction}}
\UnaryInf$ \fCenter !A \lor \lnot !A$
\end{prooftree}
\end{ex}

\begin{prob}
Maak !!{derivation}s van de volgende sequenten:
\begin{enumerate}
\item $!A \land (!B \land !C) \Sequent (!A \land !B) \land !C$.
\item $!A \lor (!B \lor !C) \Sequent (!A \lor !B) \lor !C$.
\item $!A \lif (!B \lif !C) \Sequent !B \lif (!A \lif !C)$.
\item $!A \Sequent \lnot\lnot !A$.
\end{enumerate}
\end{prob}

\begin{prob}
Maak !!{derivation}s van de volgende sequenten:
\begin{enumerate}
\item $(!A \lor !B) \lif !C \Sequent !A \lif !C$.
\item $(!A \lif !C) \land (!B \lif !C) \Sequent (!A \lor !B) \lif !C$.
\item $\Sequent \lnot(!A \land \lnot !A)$.
\item $!B \lif !A \Sequent \lnot !A \lif \lnot !B$.
\item $\Sequent (!A \lif \lnot !A) \lif \lnot !A$.
\item $\Sequent \lnot(!A \lif !B) \lif \lnot !B$.
\item $!A \lif !C \Sequent \lnot (!A \land \lnot !C)$.
\item $!A \land \lnot !C \Sequent \lnot (!A \lif !C)$.
\item $!A \lor !B, \lnot !B \Sequent !A$.
\item $\lnot !A \lor \lnot !B \Sequent \lnot(!A \land !B)$.
\item $\Sequent (\lnot !A \land \lnot !B) \lif\lnot(!A \lor !B)$.
\item $\Sequent \lnot(!A \lor !B) \lif (\lnot !A \land \lnot !B)$.
\end{enumerate}
\end{prob}

\begin{prob}
Maak !!{derivation}s van de volgende sequenten:
\begin{enumerate}
\item $\lnot(!A \lif !B) \Sequent !A$.
\item $\lnot(!A \land !B) \Sequent \lnot !A \lor \lnot !B$.
\item $!A \lif !B \Sequent \lnot !A \lor !B$.
\item $\Sequent \lnot \lnot !A \lif !A$.
\item $!A \lif !B, \lnot !A \lif !B \Sequent !B$.
\item $(!A \land !B) \lif !C \Sequent (!A \lif !C) \lor (!B \lif !C)$.
\item $(!A \lif !B) \lif !A \Sequent !A$.
\item $\Sequent (!A \lif !B) \lor (!B \lif !C)$.
\end{enumerate}
(Voor al deze opgaven is de regel $\RightR{\Contraction}$ nodig.)
\end{prob}
\end{document}
